\bookchapter{บทที่ 40: Requirement Engineering — เปลี่ยนไอเดียให้เป็นคำสั่งที่ AI ทำตามได้}

\chapterintro{ถ้า Context Engineering แก้ปัญหาเรื่อง “ทำให้ AI รู้เบื้องหลัง” แล้ว Requirement Engineering ก็แก้ปัญหาเรื่อง “ทำให้ AI รู้ว่าครั้งนี้ต้องทำอะไรกันแน่” นี่คือขั้นตอนที่เผยความต่างระหว่างคนกับคนได้ชัดที่สุดในการเขียนโค้ดด้วย AI — เครื่องมือเหมือนกัน โมเดลเหมือนกัน แต่บางคนสร้างโปรเจกต์ที่น่าทึ่งออกมาได้ บางคนก็ได้แค่ของเล่นธรรมดา ความต่างมักไม่ได้อยู่ที่ฝีมือทางเทคนิค แต่อยู่ที่ความสามารถในการสื่อสารความต้องการ}

\booksection{จากไอเดียคลุมเครือสู่ความต้องการที่มีโครงสร้าง}

ทุกโปรเจกต์เริ่มต้นจากไอเดียคลุมเครือ — “อยากทำแอปบันทึกรายรับรายจ่าย” “อยากทำเว็บส่วนตัว” “อยากทำเครื่องมือ to-do list” ไอเดียแบบนี้ไม่มีอะไรผิด แต่ถ้าคุณโยนมันให้ AI ไปตรง ๆ คุณภาพผลงานจะขึ้นอยู่กับดวงล้วน ๆ เป้าหมายของ Requirement Engineering คือการเปลี่ยนไอเดียคลุมเครือแบบนี้ ให้กลายเป็นความต้องการที่มีโครงสร้างซึ่ง AI ทำตามได้อย่างแม่นยำ

ข่าวดีคือ: ความต้องการแบบมีโครงสร้างไม่ได้ต้องการให้คุณเขียน PRD สักเล่มหนา แค่ตอบคำถามห้าข้อก็พอ ซึ่งเราเรียกว่า “องค์ประกอบ 5 ข้อของความต้องการ”:

\begin{center}
\begin{tabularx}{\linewidth}{|X|X|X|}
\hline
\textbf{องค์ประกอบ} & \textbf{คำถาม} & \textbf{ตัวอย่าง: แอปบันทึกรายรับรายจ่าย} \tabularnewline
\hline
Who & ให้ใครใช้? & ใช้ส่วนตัว บันทึกรายรับรายจ่ายประจำวัน \tabularnewline
What & ทำอะไร? & แอปที่บันทึกรายรับรายจ่ายทีละรายการ สรุปเป็นรายเดือน และแสดงผลเป็นภาพได้ \tabularnewline
Why & ทำเพื่ออะไร? & อยากเข้าใจโครงสร้างการใช้จ่ายของตัวเอง เพื่อหาค่าใช้จ่ายที่ไม่จำเป็น \tabularnewline
How & ทำอย่างไรโดยคร่าว ๆ? & ฝั่ง front-end ใช้ React เก็บข้อมูลไว้ในเบราว์เซอร์ก่อน ยังไม่ทำ back-end และระบบล็อกอิน \tabularnewline
Constraint & มีข้อจำกัดอะไร? & พัฒนาคนเดียว เสร็จภายใน 2 สัปดาห์ ไม่ต้อง deploy ขึ้นใช้งานจริง \tabularnewline
\hline
\end{tabularx}
\end{center}

\begin{center}\small\color{gray}องค์ประกอบ 5 ข้อของความต้องการพร้อมตัวอย่าง\end{center}

องค์ประกอบทั้งห้าไม่จำเป็นต้องเขียนเป็นเอกสารทางการ แม้แต่แค่ไล่คิดในหัวให้จบเร็ว ๆ ก็เพียงพอแล้ว กุญแจสำคัญคือสร้างนิสัย “คิดให้ชัดก่อนเปิดปาก” คุณจะพบว่าพอตอบห้าคำถามนี้ได้ชัดเจน คำอธิบายความต้องการของคุณก็เป็นแบบมีโครงสร้างไปเอง

\booksection{ศิลปะการแบ่งชั้นของการแตกความต้องการ}

สำหรับโปรเจกต์ที่ใหญ่ขึ้น คำอธิบายความต้องการชุดเดียวยังไม่พอ คุณต้องแตกความต้องการใหญ่ออกเป็นความต้องการย่อยที่เล็กลงและตรวจสอบได้ด้วยตัวเอง กระบวนการนี้เรียกว่า “การแตกความต้องการ” หลักการสำคัญคือ: ความต้องการย่อยแต่ละชิ้นต้องทำเสร็จได้อย่างอิสระ ตรวจสอบได้อย่างอิสระ และส่งมอบได้อย่างอิสระ

ยกตัวอย่าง “ทำแอปบันทึกรายรับรายจ่ายส่วนตัว” วิธีแตกงานที่สมเหตุสมผลคือ:

\begin{enumerate}
  \item ทำฟังก์ชันบันทึกรายรับรายจ่ายพื้นฐาน (เพิ่ม ดู และลบรายการได้);
  \item เพิ่มฟังก์ชันหมวดหมู่ (ติดป้ายกำกับแต่ละรายการว่าอาหาร การเดินทาง ฯลฯ);
  \item ทำฟังก์ชันสรุปรายเดือน (สรุปภาพรวมรายรับรายจ่ายเป็นรายเดือน);
  \item เพิ่มการแสดงข้อมูลเป็นภาพ (ใช้กราฟโชว์โครงสร้างการใช้จ่าย)
\end{enumerate}

แต่ละขั้นตอนคือฟังก์ชันที่เพิ่มเข้ามาทีละส่วน และรันทดสอบได้ด้วยตัวเอง

กฎทองของการแตกโจทย์: จบแต่ละขั้น โปรเจกต์ต้องเป็น "ของที่ใช้ได้" แม้ฟีเจอร์จะยังหยาบ อย่าปล่อยให้โปรเจกต์ค้างอยู่ในสภาพ "ทำไปครึ่งทางทุกอย่างแต่ใช้ไม่ได้สักอย่าง"

การแตกงานแบบนี้ได้ประโยชน์สองอย่าง หนึ่ง มันทำให้คุณเห็นความคืบหน้าได้ตลอด — ทำเสร็จทีละขั้น แอปของคุณก็ได้ฟังก์ชันเพิ่มขึ้นหนึ่ง ผลตอบรับเชิงบวกแบบนี้สำคัญมากสำหรับโปรเจกต์ระยะยาว (คนเราก็ต้องการความรู้สึกสำเร็จนี่นา) สอง มันทำให้งานของ AI แต่ละครั้งเล็กพอดีตัว ลดความเสี่ยงเรื่อง “AI เข้าใจคลาดเคลื่อน” — งานยิ่งเล็ก AI ยิ่งเข้าใจแม่น คุณภาพผลงานยิ่งสูง

\booksection{เทมเพลตการเขียน PRD และการใช้งานจริง}

สำหรับโปรเจกต์ที่ใหญ่ขึ้น (เช่น แอป full-stack ที่มีทั้ง front-end, back-end และโมดูลฟังก์ชันหลายตัว) ใช้แค่องค์ประกอบ 5 ข้ออาจไม่พอ ถึงจุดนั้นคุณต้องเขียนเอกสารความต้องการโปรเจกต์ (PRD) ที่เป็นทางการขึ้นมาหนึ่งชุด จุดที่มักเข้าใจผิดคือ: PRD ไม่ได้เขียนให้ AI อ่าน — มันเขียนให้ตัวคุณเอง ช่วยจัดระเบียบความคิดให้คุณ จากนั้นคุณถึงจะสั่งงาน AI ได้แม่นยำขึ้น

PRD ที่ดีไม่จำเป็นต้องยาว แต่ควรมีหกส่วน: ภาพรวมโปรเจกต์ (บอกให้ชัดในหนึ่งย่อหน้าว่าทำอะไร) รายการฟังก์ชัน (แต่ละฟังก์ชันคืออะไร) แผนเทคนิค (ใช้ tech stack อะไร สถาปัตยกรรมโดยรวมเป็นแบบไหน) โมเดลข้อมูล (โครงสร้างข้อมูลหลัก) คำอธิบายหน้าจอ (หน้าหลัก ๆ หน้าตาเป็นอย่างไร) และแผนการพัฒนา (แบ่งทำกี่ขั้น)

\begin{lstlisting}
# PRD แอปบันทึกรายรับรายจ่ายส่วนตัว
## ภาพรวมโปรเจกต์
แอปบันทึกรายรับรายจ่ายส่วนตัวแบบเบา ๆ รองรับการบันทึก หมวดหมู่ สรุปรายเดือน และกราฟข้อมูล
## รายการฟีเจอร์
1. บันทึกรายรับรายจ่าย: จำนวนเงิน ประเภท (รายรับ/รายจ่าย) หมวดหมู่ โน้ต วันที่
2. จัดการหมวดหมู่: เตรียมอาหาร เดินทาง ช้อปปิ้ง บันเทิง อื่น ๆ ให้ก่อน
3. สรุปรายเดือน: รวมรายรับ รายจ่าย และยอดคงเหลือรายเดือน
4. กราฟการใช้จ่าย: วงกลมแสดงสัดส่วนแต่ละหมวด
5. ส่งออกข้อมูล: เป็นไฟล์ CSV
## แนวทางเทคนิค
- Frontend: React + TypeScript + Tailwind CSS
- เก็บข้อมูล: localStorage (ไม่ต้องมีแบ็กเอนด์)
- กราฟ: ใช้ Chart.js
- บิลด์: Vite
## โมเดลข้อมูล
Record {
  id: string
  amount: number
  type: 'income' | 'expense'
  category: string
  note: string
  date: string  // YYYY-MM-DD
}
## แผนพัฒนา
Step 1: ฟังก์ชันบันทึกพื้นฐาน (CRUD)
Step 2: จัดการหมวดหมู่
Step 3: สรุปรายเดือน
Step 4: แสดงข้อมูลเป็นกราฟ
Step 5: ส่งออกข้อมูล
\end{lstlisting}

สังเกตสไตล์การเขียน PRD แบบนี้: ใช้ลิสต์กระชับกับฟอร์แมตที่มีโครงสร้าง ไม่ใช่เรียงความยาวเหยียด แต่ละฟังก์ชันมีนิยามชัดเจน โมเดลข้อมูลอธิบายด้วย type แบบง่าย ๆ PRD สไตล์นี้ไม่ใช่แค่คุณอ่านรู้เรื่อง เมื่อส่งให้ AI แล้ว AI ยังเข้าใจแต่ละจุดความต้องการได้อย่างแม่นยำด้วย

\booksection{ลงสนามจริง: ไล่ครบวงจรจากไอเดียสู่แผนงาน}

ทฤษฎีล้วน ๆ มันแห้งไป เรามาไล่กระบวนการ Requirement Engineering ทั้งหมดด้วยตัวอย่างเต็ม ๆ สักรอบ สมมุติว่าไอเดียของคุณคือ: “อยากทำ learning tracker ที่จดให้ได้ว่าแต่ละวันเรียนอะไรไป”

ขั้นแรก ไล่องค์ประกอบ 5 ข้อของความต้องการ

\begin{itemize}
  \item Who: ใช้ส่วนตัว
  \item What: เครื่องมือที่จดเนื้อหาที่เรียนแต่ละวัน ติดตามเวลาเรียน และทบทวนความคืบหน้าได้
  \item Why: อยากสร้างนิสัยเรียนรู้ต่อเนื่อง และได้เห็นการสะสมของตัวเอง
  \item How: เริ่มจากทำเวอร์ชันเว็บ เก็บข้อมูลไว้ในเบราว์เซอร์
  \item Constraint: พัฒนาโดยคนที่เริ่มจากศูนย์ เสร็จฟังก์ชันหลักภายใน 1 สัปดาห์
\end{itemize}

ขั้นที่สอง แตกความต้องการ

\begin{enumerate}
  \item Step 1: ทำฟังก์ชัน "บันทึกรายการการเรียน" (วันที่ เนื้อหา ระยะเวลา);
  \item Step 2: เพิ่ม "สถิติเวลาเรียน" (วันนี้ และสะสมทั้งสัปดาห์);
  \item Step 3: ทำมุมมอง "ปฏิทินการเรียน" (แสดงแบบกราฟ contribution ของ GitHub);
  \item Step 4: เพิ่มฟังก์ชัน "ทบทวนการเรียน" (สุ่มโชว์บันทึกการเรียนของวันที่ผ่านมา ช่วยให้ทบทวนได้)
\end{enumerate}

ขั้นที่สาม เขียน PRD แบบสั้น ไม่ต้องยาว — ผลการแตกงานข้างบนเองก็เป็นเนื้อหาหลักของ PRD แล้ว เอาไปจัดเป็นไฟล์ .md หนึ่งไฟล์ วางไว้ที่รากของโปรเจกต์ ใช้เป็นข้อมูลอ้างอิงร่วมกันระหว่างคุณกับ AI

ขั้นที่สี่ เริ่มพัฒนาขั้นตอนแรก คำสั่งที่ให้ AI อาจเป็นแบบนี้:

\begin{lstlisting}
ต้องการฟีเจอร์แรกของตัวช่วยติดตามการเรียน—บันทึกรายการเรียนรู้
ผู้ใช้เพิ่มบันทึกได้หนึ่งรายการ ประกอบด้วยวันที่ เนื้อหาที่เรียน และระยะเวลา (นาที)
เก็บข้อมูลใน localStorage
ฝั่งหน้าเว็บใช้ React + TypeScript
ขอโค้ดเต็มมาให้หน่อย
\end{lstlisting}

สังเกตจุดเด่นของคำสั่งนี้: มันอธิบายเป้าหมายของขั้นตอนปัจจุบันเท่านั้น (ฟังก์ชันบันทึก) ไม่แทรกฟังก์ชันอื่น; มันระบุข้อจำกัดทางเทคนิคชัดเจน (React, TypeScript, localStorage); มันให้ข้อมูลมากพอให้ AI รู้ว่า “ต้องทำอะไร” และ “ทำอย่างไร” แต่วิธีลงมือจริงให้ AI เป็นคนตัดสิน นี่คือคำสั่งเขียนโค้ดด้วย AI คุณภาพสูงที่ผ่านการขัดเกลาด้วย Requirement Engineering มาแล้ว

\begin{tcbpractice}{ลงมือทำ: เขียนบัตรองค์ประกอบ 5 ข้อให้ "โปรเจกต์ในใจ" ของคุณ}
เลือกโปรเจกต์เล็ก ๆ ที่คุณหมักไว้นานแล้วสักหนึ่ง ใช้เวลา 5 นาทีกรอกบัตรหนึ่งใบ: Who / What / Why / How / Constraint อย่างละหนึ่งประโยค จากนั้นแตกมันออกเป็น 3-5 ขั้นตอนแบบ “จบแต่ละขั้นแล้วใช้งานได้” คุณจะพบว่า: พอเขียนไอเดียลงกระดาษ มันก็เปลี่ยนจาก “ฝันกลางวัน” มาเป็น “to-do list” เสียแล้ว

\end{tcbpractice}

\begin{tcbsummary}{สรุปบทนี้}
\begin{itemize}
  \item องค์ประกอบ 5 ข้อของความต้องการ: Who / What / Why / How / Constraint — คิดให้ชัดก่อนเปิดปาก
  \item กฎทองของการแตกความต้องการ: จบแต่ละขั้นตอนแล้ว โปรเจกต์ต้องเป็น "ของที่ใช้ได้" เสมอ
  \item PRD เขียนให้ตัวเอง: ภาพรวม รายการฟังก์ชัน แผนเทคนิค โมเดลข้อมูล หน้าจอ และแผนงาน
  \item คำสั่งที่ดี = อธิบายแค่ขั้นตอนปัจจุบัน + ระบุข้อจำกัดชัด + บอก "ทำอะไร" ไม่ใช่ "ทำอย่างไร"
  \item เทมเพลตไล่ครบวงจร: 5 องค์ประกอบ แตกงาน PRD สั่งงาน AI ครบในสี่ขั้น
\end{itemize}

\end{tcbsummary}
